

\documentclass[12pt]{article}

\usepackage[margin = .8in]{geometry}
\usepackage{amsmath}
\usepackage{graphicx}
\usepackage{multicol, enumerate, tabularx}

\usepackage[parfill]{parskip}

\usepackage{adjustbox, soul}

\usepackage{fancyhdr}
\pagestyle{fancy}

\lhead{\emph{Math F113X: Substitution and Transposition Ciphers}}
\rhead{\emph{lecture notes}}

\usepackage{tikz}
\usetikzlibrary{calc,trees,positioning,arrows,fit,shapes,through, backgrounds}
\usetikzlibrary{patterns}

\usetikzlibrary{decorations.markings}
\usetikzlibrary{arrows}

\renewcommand{\emph}[1]{\textsf{\textbf{#1}}}
%\newcommand{\ans}[1][1in]{\rule{#1}{.5pt}}

\usepackage{pgfplots}

\usepackage{longtable}
\usepackage{tabularx}

\newcommand{\ds}{\displaystyle}
\newcommand{\ans}[1][1in]{\rule{#1}{.5pt}}

\newcommand{\points}[1]{(#1 points.)}		% Trying to be lazy.

\usepackage{array}
\newcolumntype{L}[1]{>{\raggedright\let\newline\\\arraybackslash\hspace{0pt}}m{#1}}
\newcolumntype{C}[1]{>{\centering\let\newline\\\arraybackslash\hspace{0pt}}m{#1}}
\newcolumntype{R}[1]{>{\raggedleft\let\newline\\\arraybackslash\hspace{0pt}}m{#1}}
\newcommand{\red}[1]{\textcolor{red}{#1}}

\newcommand{\be}{\begin{enumerate}}
\newcommand{\ee}{\end{enumerate}}

%\topmargin -1in
%\textheight 9.5in
%\oddsidemargin -0.3in
%\evensidemargin \oddsidemargin
%\pagestyle{empty}
%%\marginparwidth 0.5in
%\textwidth 7in
%\parindent 0in

%--------------------------------------------------------------------------------------------------------------------------------------------------------------------------
%						Document
%--------------------------------------------------------------------------------------------------------------------------------------------------------------------------


\begin{document}
%\pagestyle{fancy}
%\begin{center}
%{\Large  Intro Cryptography (Day 2)}
%\end{center}
\begin{enumerate}
\item Review \emph{shift ciphers}: What is the key in the shift cipher below? KEY: \ans\\

\hspace*{-2.2cm}
{\scriptsize
\begin{tabular}{|c|c|c|c|c|c|c|c|c|c|c|c|c|c|c|c|c|c|c|c|c|c|c|c|c|c|c|}
\hline
&0&1&2&3&4&5&6&7&8&9&10&11&12&13&14&15&16&17&18&19&20&21&22&23&24&25\\ \hline \hline
in& A&B&C&D&E&F&G&H&I&J&K&L&M&N&O&P&Q&R&S&T&U&V&W&X&Y&Z\\  \hline
out&L &M&N&O&P&Q&R&S&T&U&V&W&X&Y&Z&A&B&C&D&E&F&G&H&I&J&K\\[12pt]
\hline
\end{tabular}
}
\item A shift cipher is a particular type of \emph{substitution cipher}. Below is a substitution cipher that is \textbf{not} a shift cipher.\\


{\scriptsize
\begin{tabular}{|c|c|c|c|c|c|c|c|c|c|c|c|c|c|c|c|c|c|c|}
\hline
&0&1&2&3&4&5&6&7&8&9&10&11&12&13&14&15&16&17\\ \hline \hline
in& A&B&C&D&E&F&G&H&I&J&K&L&M&N&O&P&Q&R\\  \hline
out&E&F&G&L&8&A&R&Q&T&U&V&P&B&D&N&O&H&M\\
\hline \hline
&18&19&20&21&22&23&24&25&26&27&28&29&30&31&32&33&34&35\\ \hline \hline
in&S&T&U&V&W&X&Y&Z&0&1&2&3&4&5&6&7&8&9\\ \hline
out&I&1&4&7&J&6&5&K&9&3&2&0&C&X&S&Z&Y&W \\
[12pt]
\hline
\end{tabular}
}
	\begin{enumerate}
	\item Encrypt: \textsf{MEET APRIL 3}
	
	\vfill
	
	\item Decrypt: \textsf{F 4 I E 1 9 Z 0 9}
	
	\vfill
	
	\item What key is needed to decrypt a message using this encryption scheme?
	
	\vfill
	
	\item Which substitution is, in general, harder to break, a shift cipher or one that is not a shift cipher?
	
	\vfill
	
	\item What strategies would you use to try to break a substitution cipher that is not a shift cipher?
	\vfill
	\end{enumerate}
	
	\newpage
	
\item Another encryption scheme is called a \emph{transposition cipher}. 

Encode \textsf{JEWEL FOUND} using a transposition cipher of rows of 4 letters (4 columns).

\be
\item Put in rows
	
\item read out columns

\ee

\vfill

\item Decrypt \begin{quote} \textsf{BNEAA \ ANBDA \ RX } \end{quote}
if it was encoded with a transposition cipher with  rows of 4 letters (4 columns).

\be
\item How many rows? 

\item Put in columns
	
\item read out rows

\ee

\vfill

\item Encrypt using a transposition cipher with a keyword {\sf KEY}: \begin{quote}\textsf{ H E L P M E }\end{quote}
\be
\item Put in grid by rows

\item Mix up columns

\item Read out columns

\ee

\vfill

\item Decrypt using a transposition cipher with a keyword {\sf KEY}: 
\begin{quote} \textsf{F N T I U \ \ I O D X} \end{quote}

\be
\item Determine how many rows

\item Put in grid by \emph{columns}, labeled by SORTED keyword

\item un-mix-up columns

\item Read out \emph{rows}

\ee

\vfill


\end{enumerate}
\end{document}

%-------------------------------------------------------------------------------------------------------------------------------------------------------------------------------------------------------------------

%%% Local Variables:
%%% mode: latex
%%% TeX-master: t
%%% End:
